% !TeX root = ../se200_identity_regimes.tex

\section{Limits and Boundary Conditions}
\label{sec:limits}

This section specifies the boundary of applicability of the
result and identifies classes of systems
excluded by the admissibility constraints.

\subsection{Scope of the Result}
\label{subsec:limits_scope}

The result applies only to admissible substrates
(Section~\ref{sec:preliminaries}),
namely those in which identity conditions, persistence conditions,
and admissible relations are invariant under admissible extension
and independent of interpretive commitment.

In particular, the result concerns neutral substrates under
persistent interpretive disagreement, and establishes only a
necessary lower bound on the number of regime profiles required
to maintain stable identity under such conditions.

The result is structural:
it characterizes identity and persistence behavior
in terms of invariance under a fixed transformation basis,
and does not depend on any domain-specific interpretation.

\subsection{Excluded Systems}
\label{subsec:limits_excluded}

The following classes of systems are excluded from the domain of the result:

\begin{itemize}

  \item \textbf{Interpretively dependent identity:}
        Systems in which the identity condition of an entity
        depends on interpretive commitments realized in extensions.
        Such systems violate identity invariance and
        do not admit a neutral substrate.

  \item \textbf{Relation-determined identity:}
      Systems in which relations introduce or alter identity conditions
      for their relata.
      Relations defined over pre-established
      equivalence classes are admissible;
      relations that constitute or modify those classes are not.

  \item \textbf{Extension-variant identity or relations:}
        Systems in which identity conditions or admissible relations
        vary under extension of the represented domain.
        Such systems violate extension stability.

  \item \textbf{Embedded interpretive commitments:}
        Systems that encode interpretive content at the substrate level
        whose validity depends on extension-specific assumptions.
        Such systems violate interpretive non-commitment.

  \item \textbf{Non-admissible transformation semantics:}
        Systems requiring transformation behavior not representable
        within the fixed transformation basis $\mathbb{F}$.
        Such systems fall outside the structural scope of the analysis.

\end{itemize}

\subsection{Non-Goals}
\label{subsec:limits_nongoals}

This paper does not:

\begin{itemize}
  \item establish an upper bound or claim completeness of the regime set;
  \item construct or analyze interpretive extensions;
  \item provide domain-specific modeling guidance;
  \item assign causal or normative interpretation to structural representations;
  \item establish completeness or minimality of the transformation basis
      $\mathbb{F}$, or determine whether additional transformation families
      would force further regime splits.
\end{itemize}

\subsection{Boundary of Representation}
\label{subsec:limits_boundary_representation}

The representation preserves identity, regime typing,
and admissible relations under transformation.
It does not preserve interpretive content.

Accordingly:

\begin{itemize}
\item causal and normative interpretations arise only in extensions,
\item interpretive disagreement does not affect identity or persistence,
\item structural equivalence does not imply interpretive equivalence.
\end{itemize}


The limits identified in this section are consequences
of the neutrality constraint and the restriction
to admissible transformation behavior.
They are not limitations of a particular representation technique,
but intrinsic to the structural setting considered.

